% !TEX root = ../Projektdokumentation.tex
\section{Analysephase}
\label{sec:Analysephase}


\subsection{Ist-Analyse}
\label{sec:IstAnalyse}
Der \ac{PVE}-Cluster der \ac{ODCF} besteht aus insgesamt sechs
Hypervisor-Knoten -- drei älteren IBM- bzw.\ Lenovo-Servern und drei
neueren Dell-VxRail-Knoten -- und betreibt derzeit etwa 60 \acp{VM}
sowie eine Reihe von \acp{LXC} mit Forschungs- und
Infrastrukturdiensten.
Als primäres Storage-Backend wird ein \ac{Ceph}-Cluster mit einer
Brutto-Kapazität von rund 44\,TB eingesetzt, von denen aktuell ca.\
6\,TB belegt sind.

Die bestehende Sicherung ist heterogen aufgebaut:
\begin{itemize}
    \item Sicherungen der \acp{VM} und Container werden auf den
          \emph{lokalen} Storages der älteren Hypervisor-Knoten
          abgelegt. Die neueren Dell-VxRail-Knoten besitzen
          \textbf{keinen} lokalen Storage; für die dort laufenden
          \acp{VM} gibt es daher keine lokale Sicherungsmöglichkeit.
    \item Ausgewählte \acp{VM} werden zusätzlich per \texttt{rsync} auf
          einen separaten Backup-Server kopiert.
    \item Vom Backup-Server werden die Daten über \ac{TSM} auf Tape
          ausgelagert.
\end{itemize}

\subsubsection*{Schmerzpunkte der heutigen Lösung}
\begin{itemize}
    \item \textbf{Restore-Reichweite:} Restores aus dem lokalen
          Storage einzelner Hypervisor sind zwar schnell (typisch
          2--6\,min), funktionieren aber nur auf den älteren IBM-/
          Lenovo-Knoten. Für die VxRail-Hypervisor existiert keine
          äquivalente schnelle Restore-Quelle.
    \item \textbf{Fehlschläge bei Migrationen:} Bei der Migration
          großer \acp{VM} zwischen Knoten -- insbesondere von und zu
          VxRails -- kommt es vor, dass keine konsistente Sicherung
          mehr zur Verfügung steht.
    \item \textbf{Vollgelaufener Backup-Storage:} Da pro Hypervisor
          lokal gesichert wird, läuft der lokale Storage einzelner
          Knoten gelegentlich voll und führt zu fehlgeschlagenen
          Backup-Läufen.
    \item \textbf{Keine zentrale Verwaltung:} Backup-Jobs, Retention
          und Status sind über alle Hypervisor verstreut und werden
          nicht zentral überwacht.
    \item \textbf{Keine Integritätsprüfung:} Eine automatische
          Verifikation der gespeicherten Sicherungen findet nicht
          statt.
    \item \textbf{Keine Trennung Primär-/Sicherungsdaten:} Ein
          Ausfall eines Hypervisors zieht den lokalen Backup-Bestand
          desselben Hypervisors mit in Mitleidenschaft.
\end{itemize}


\subsection{Wirtschaftlichkeitsanalyse}
\label{sec:Wirtschaftlichkeitsanalyse}
Da konkrete interne Stundensätze und Hardware-Listenpreise nicht
abschließend vorliegen, wird mit Platzhalterwerten kalkuliert. Die
Projektkosten setzen sich zusammen aus ca.\ 40\,h Eigenleistung des
Prüflings ($\approx$\,1\,000\,\texteuro\ bei 25\,\texteuro/h
Vollkostensatz), ca.\ 5\,h Betreuung ($\approx$\,400\,\texteuro\ bei
80\,\texteuro/h) sowie geschätzten 200\,\texteuro\ für Strom und
Stellplatz des nachgenutzten IBM x3755 -- in Summe ca.\
1\,600\,\texteuro. Hardware- und Lizenzkosten fallen nicht an, da
vorhandene Hardware genutzt wird und der \ac{PBS} unter AGPLv3 frei
verfügbar ist.

Der Nutzen ergibt sich vor allem aus der Trennung von Primär- und
Sicherungsdaten sowie deutlich kürzeren Restore-Zeiten gegenüber
\ac{TSM}/Tape; bereits ein einziger vermiedener mehrstündiger Restore-
Vorgang pro Jahr amortisiert die Projektkosten. Hinzu kommt die
zentrale Verwaltung und automatische Integritätsprüfung, die den
laufenden administrativen Aufwand reduziert.


\subsection{Lasten-/Pflichtenheft}
\label{sec:Lastenheft}
Als Lastenheft dient der von der \ac{ODCF} freigegebene Projektantrag,
der Ziel, Umfang und zu treffende Entscheidungen festlegt.
Ein gesondertes Pflichtenheft wird nicht erstellt; die konkreten
Umsetzungsentscheidungen werden stattdessen in der Entwurfsphase
(Abschnitt~\ref{sec:Entwurfsphase}) dokumentiert.
Ein Auszug aus den zentralen Anforderungen befindet sich im
\Anhang{app:Lastenheft}.
